\documentclass[10pt,a4paper]{article}

\usepackage[margin=1.25cm]{geometry}
\usepackage{amsmath,amssymb,booktabs,multicol}
\usepackage{enumitem}
\usepackage[T1]{fontenc}
\usepackage[utf8]{inputenc}
\usepackage{lmodern}
\usepackage{microtype}

\setlength{\parindent}{0pt}
\setlength{\parskip}{2.5pt}
\setlength{\columnsep}{0.55cm}
\setlist[itemize]{leftmargin=*,nosep}
\pagestyle{empty}

\newcommand{\code}[1]{\texttt{\small #1}}
\newcommand{\qualitysuccess}{q_{\mathrm{success}}}
\newcommand{\qualityspeed}{q_{\mathrm{speed}}}

\begin{document}

\begin{center}
    {\Large\bfseries ProviderRouterPR1: Scoring and Health Mechanisms}\\[-1pt]
    {\small Concise summary of the currently implemented behaviour}
\end{center}

\vspace{-2pt}
\fbox{\parbox{0.975\linewidth}{\centering\small
Configured providers
$\rightarrow$ \textbf{hard eligibility and health filter}
$\rightarrow$ tie-break order
$\rightarrow$ recent metrics
$\rightarrow$ \textbf{score and stable sort}
$\rightarrow$ best-first attempts with fallback.
\quad Observed outcomes feed persistent metrics; provider-attributable outcomes
also update in-memory health.
}}

\vspace{3pt}
\begin{multicols}{2}
\small

\textbf{Scoring evidence.}
The router attempts to persist one \code{MetricsEvent} per observed provider
outcome; excluded providers produce none. An early-closed stream with no observed
outcome is deliberately not recorded. Aggregation is per provider and keeps
regular and streaming calls separate:
regular latency is time to the complete response, whereas streaming latency is
time to first chunk (TTFT). For the selected bucket, let $w_i$ be an event weight,
$y_i\in\{0,1\}$ its success indicator, and $\ell_i$ its latency:
\[
 A=\sum_i w_i,\qquad S=\sum_iw_i y_i,\qquad r=\frac{S}{A}\ (A>0),
\quad
\bar{\ell}=
\frac{\sum_{i\in L}w_i\ell_i}{\sum_{i\in L}w_i},
\]
where $L$ contains only successful events with a recorded latency. Failed
latencies are deliberately ignored: otherwise a provider that fails quickly
could appear fast. A missing rate or latency is represented by \code{None},
never by zero.

By default, events in the latest 336 hours have $w_i=1$. With a positive
half-life $h$, this window is replaced by the latest $6h$ and
\[
 w_i=2^{-a_i/h},
\]
where $a_i$ is event age in hours. Decay changes scoring evidence but not exact
diagnostic error, rate-limit, or timeout counts.

\textbf{Score calculation.}
The raw speed quality is a bounded inverse-latency curve
\[
 q_{\mathrm{raw}}=\frac{R}{R+\bar{\ell}},
\]
with $R=2000$\,ms for regular calls and $R=500$\,ms for streaming TTFT.
Thus latency $R$ maps to $0.5$. If no latency exists, $q_{\mathrm{raw}}=0$.

Both observed qualities are shrunk toward an optimistic prior $q_0$:
\[
 \qualitysuccess=\frac{n_0q_0+Ar}{n_0+A},
 \qquad
 \qualityspeed=\frac{n_0q_0+S q_{\mathrm{raw}}}{n_0+S}.
\]
Defaults are $q_0=0.75$ and $n_0=5$ pretend attempts. Consequently, a new
provider scores exactly $0.75$; sparse observations move it gradually, while
deep history dominates the prior. This supplies exploration without a separate,
stateful exploration bonus.

The final score is the normalized weighted average
\[
 Q=\frac{\alpha\qualitysuccess+\beta\qualityspeed}{\alpha+\beta},
\qquad \alpha,\beta\geq0,\quad \alpha+\beta>0,
\]
with $\alpha=\beta=1$ by default. Normalization keeps $Q\in[0,1]$ and makes the
weights relative rather than scale-dependent. Error counts are not separate
factors: failures already reduce $r$, so another penalty would double-count
them. Cost is not currently scored.

\textbf{Ordering and failure behaviour.}
\code{ScoreBasedPolicy} receives only eligible providers. It first invokes one
tie-break policy (round robin by default), then queries and scores history on
every call and applies one stable descending sort. Equal scores retain the
round-robin order. The full order is handed to fallback, so failure of the
highest-ranked provider tries the next one. With disabled metrics, no eligible
providers, or a failed history query, routing safely returns the tie-break order;
scores are neither cached nor allowed to break an LLM call.

\columnbreak

\textbf{Health is a hard gate, not a score.}
Health answers ``may this provider be tried now?'', while scoring answers ``which
eligible provider should be tried first?'' A benched provider is excluded before
the policy runs rather than assigned an arbitrary numerical penalty. This avoids
mixing operational safety with performance preference and lets any routing
policy reuse the same health state.

\begin{center}
\small
\begin{tabular}{@{}p{0.28\columnwidth}p{0.65\columnwidth}@{}}
\toprule
\textbf{Failure} & \textbf{Implemented transition}\\
\midrule
HTTP 401/403 &
Disable provider until explicit reset or destruction of the router instance.\\
HTTP 429 &
Immediate cooldown (default 60\,s); do not change the consecutive-failure count.\\
Timeout/408, connection, ${\geq}500$, interrupted stream, unknown &
Increment consecutive failures; at the default threshold of 3, start a
60\,s cooldown.\\
HTTP 400/422 or invalid operation/arguments &
Stop the call because another provider using the same request would not help;
do not blame or bench the provider.\\
Success &
Clear failure count, cooldown, trigger, and last error; an authentication bench
remains until explicitly reset.\\
\bottomrule
\end{tabular}
\end{center}

Cooldowns use a monotonic clock. Expiry restores eligibility but deliberately
does not erase the failure count; a failed probe can therefore re-bench the
provider immediately. \code{reset\_health()} restores eligibility but never
deletes metrics. This separates short-lived, process-local protection from
longer-lived performance evidence. A bench created during a fallback sequence
affects the next invocation because the current best-first list has already
been fixed.

\textbf{Important architectural choices.}
\begin{itemize}
    \item Aggregation and scoring are pure Python functions, independent of
    adapters and storage backends; the policy only composes them.
    \item Successful latency alone represents speed; every attempt represents
    reliability. Regular and streaming evidence is never blended.
    \item \code{use\_streaming} is fixed when constructing the policy, not inferred
    from opaque provider arguments. Mixed workloads therefore require separate
    policy instances or an explicit choice of which history to use.
    \item Provider names are metric and health identity keys. Names must be
    unique, and renaming a provider intentionally starts a new effective history.
    \item Metric writes and reads are best-effort. Storage failure may reduce
    future evidence, but never replaces a successful model response.
\end{itemize}

\textbf{Short planned extensions.}
\begin{itemize}
    \item Optional user-configured token-cost calculation; cost is deliberately
    outside the current core score.
    \item Pre-flight capability filtering inferred from call arguments; currently
    unsupported capabilities fail during the provider call.
\end{itemize}

\textbf{Implementation basis.}
\code{scoring.py}, \code{stats.py}, \code{policies/score\_based.py},
\code{health.py}, \code{filters.py}, \code{metrics.py}, and \code{router.py}.

\end{multicols}
\end{document}
